\documentclass[xetex,mathserif,serif]{beamer}

\usepackage{xunicode}
\usepackage{xltxtra}
\usepackage{color}
\usepackage{url}
\usepackage{listings}
\usepackage{fontspec}
\usepackage{geometry}
\usepackage{lastpage}
\usepackage{fancyhdr}
\usepackage{amsmath}
\usepackage{amsthm}
\usepackage{amssymb}
\usepackage{blkarray}
\usepackage{multicol}
\usepackage{relsize}

\definecolor{solarized@base03}{HTML}{002B36}
\definecolor{solarized@base02}{HTML}{073642}
\definecolor{solarized@base01}{HTML}{586e75}
\definecolor{solarized@base00}{HTML}{657b83}
\definecolor{solarized@base0}{HTML}{839496}
\definecolor{solarized@base1}{HTML}{93a1a1}
\definecolor{solarized@base2}{HTML}{EEE8D5}
\definecolor{solarized@base3}{HTML}{FDF6E3}
\definecolor{solarized@yellow}{HTML}{B58900}
\definecolor{solarized@orange}{HTML}{CB4B16}
\definecolor{solarized@red}{HTML}{DC322F}
\definecolor{solarized@magenta}{HTML}{D33682}
\definecolor{solarized@violet}{HTML}{6C71C4}
\definecolor{solarized@blue}{HTML}{268BD2}
\definecolor{solarized@cyan}{HTML}{2AA198}
\definecolor{solarized@green}{HTML}{859900}
\definecolor{yaleblue}{HTML}{0E4C92}

\setbeamertemplate{navigation symbols}{}
% \setbeamerfont{title}{family=\old}
% \setbeamerfont{author}{family=\tfont}%
% \setbeamerfont{frametitle}{family=\oldA}
% \setbeamerfont{date}{family=\dfont}

\setbeamertemplate{itemize items}{--}
\setbeamercolor*{item}{fg=black}

\defaultfontfeatures{Mapping=tex-text}
\hypersetup{pdfstartview={FitH}}

\newcommand{\old}[1]{\fontspec[Alternate=1,Ligatures={Common}]{Hoefler Text}\fontsize{18pt}{30pt}\selectfont #1}%
\newcommand{\oldA}[1]{\fontspec[Alternate=1,Ligatures={Common, Rare}]{Hoefler Text}\fontsize{12pt}{15pt}\selectfont #1}%
\newcommand{\oldB}[1]{\fontspec[Ligatures={Common}]{Didot}\fontsize{12pt}{15pt}\color{solarized@base02}\selectfont #1}%
\newcommand{\tfont}[1]{\fontspec[Alternate=1,Ligatures={Common}]{Hoefler Text}\fontsize{12pt}{20pt}\selectfont #1}%
\newcommand{\dfont}[1]{\fontspec[Ligatures={Common}]{Didot}\fontsize{12pt}{12pt}\selectfont #1}%

\newcommand{\minimize}{\mathop{\mathrm{minimize}}}
\newcommand{\argmin}{\mathop{\mathrm{arg\,min}}}
\newcommand{\argmax}{\mathop{\mathrm{arg\,max}}}
\newcommand{\st}{\mathop{\mathrm{subject\,\,to}}}

\newcommand\independent{\protect\mathpalette{\protect\independenT}{\perp}}
\def\independenT#1#2{\mathrel{\rlap{$#1#2$}\mkern2mu{#1#2}}}

\setlength{\parindent}{0pt}
\setlength{\parskip}{12pt}

\setromanfont [Ligatures={Common}, Numbers={OldStyle}, Variant=01,
 BoldFont={LinLibertine_RB.otf},
 ItalicFont={LinLibertine_RI.otf},
 BoldItalicFont={LinLibertine_RBI.otf}
 ]{LinLibertine_R.otf}



\begin{document}

%%%%%%%%%%%%%%%%%%%%%%%%%%%%%%%%%%%%%%%%%%%%%%%%%%%
\begin{frame}[fragile] \frametitle{}

\vfill

{\fontsize{0.7cm}{0cm}\selectfont Lecture 04 \\\vspace{0.2cm} Applications
and Intro to Multivariate Regression}\\\vspace{0.5cm}
14 September 2015

\vspace{2cm}

\begin{minipage}{0.6\textwidth}
Taylor B. Arnold \\
Yale Statistics \\
STAT 312/612
\end{minipage}
\hfill
\begin{minipage}{0.3\textwidth}\raggedleft
\includegraphics[scale=0.3]{../yale-logo.png}
\end{minipage}%

\end{frame}

%%%%%%%%%%%%%%%%%%%%%%%%%%%%%%%%%%%%%%%%%%%%%%%%%%%
\begin{frame}[fragile] \frametitle{}

{\color{yaleblue}\fontsize{16pt}{20pt}\selectfont Notes}

\begin{enumerate}
\item Problem set 1 due start of next class
\item TA session tomorrow night
\item One typo; 4(b) has hypothesis test for the intercept not the slope
\item Try to get a fresh copy of notes!
\end{enumerate}

\end{frame}

%%%%%%%%%%%%%%%%%%%%%%%%%%%%%%%%%%%%%%%%%%%%%%%%%%%
\begin{frame}[fragile] \frametitle{}

{\color{yaleblue}\fontsize{16pt}{20pt}\selectfont Goals for today}

\begin{enumerate}
\item Galton's heights data
\item Multivariate regression; normal equations
\item Model frames in R
\end{enumerate}

\end{frame}

%%%%%%%%%%%%%%%%%%%%%%%%%%%%%%%%%%%%%%%%%%%%%%%%%%%
\begin{frame}[fragile] \frametitle{}

\begin{flushright}
{\color{yaleblue}\sc\fontsize{1cm}{0cm}\selectfont Galton Heights Application}
\end{flushright}

\end{frame}

%%%%%%%%%%%%%%%%%%%%%%%%%%%%%%%%%%%%%%%%%%%%%%%%%%%
\begin{frame}[fragile] \frametitle{}

\begin{flushright}
{\color{yaleblue}\sc\fontsize{1cm}{0cm}\selectfont Multivariate Regression Models}
\end{flushright}

\end{frame}

%%%%%%%%%%%%%%%%%%%%%%%%%%%%%%%%%%%%%%%%%%%%%%%%%%%
\begin{frame}[fragile] \frametitle{}

The multivariate linear regression model is, on the surface,
only a slight generalization of the simple linear regression model:
\begin{align*}
y_i &= x_{1,i} \beta_1 + x_{2,i} \beta_2 + \cdots + x_{1,p} \beta_p + \epsilon
\end{align*}
\pause The statistical estimation problem now becomes one of
estimating the $p$ components of the multivariate vector $\beta$.

\end{frame}


%%%%%%%%%%%%%%%%%%%%%%%%%%%%%%%%%%%%%%%%%%%%%%%%%%%
\begin{frame}[fragile] \frametitle{}

A sample can be re-written in terms of the vector $x_i$
(the vector of covariates for a single observation):
\begin{align*}
y_i &= x_{i}^t \beta + \epsilon
\end{align*}

\end{frame}

%%%%%%%%%%%%%%%%%%%%%%%%%%%%%%%%%%%%%%%%%%%%%%%%%%%
\begin{frame}[fragile] \frametitle{}

In matrix notation, we can write the linear model simultaneously
for all observations:
\begin{align*}
\left(\begin{array}{c}y_1\\ y_2\\ \vdots\\ y_n\end{array}\right) &=
  \left(\begin{array}{cccc}x_{1,1}&x_{2,1}&\cdots&x_{p,1}\\
                           x_{1,2}&\ddots&&x_{p,2}\\
                           \vdots&&\ddots&\vdots\\
                           x_{1,n}&x_{2,n}&\cdots&x_{p,n}\\\end{array}\right)
  \left(\begin{array}{c}\beta_1\\ \beta_2\\ \vdots\\ \beta_p\end{array}\right) +
  \left(\begin{array}{c}\epsilon_1\\ \epsilon_1\\ \vdots\\ \epsilon_n\end{array}\right)
\end{align*}
\pause Which can be compactly written as:
\begin{align*}
Y &= X \beta + \epsilon
\end{align*}
\pause {\bf Note:} we use the transpose for $x_i^t \beta$ but not for $X \beta$!

\end{frame}

%%%%%%%%%%%%%%%%%%%%%%%%%%%%%%%%%%%%%%%%%%%%%%%%%%%
\begin{frame}[fragile] \frametitle{}

For reference, note the following equation
\begin{align*}
Y &= X \beta + \epsilon
\end{align*}
Yields these dimensions:
\begin{align*}
Y &\in \mathbb{R}^n \\
X &\in \mathbb{R}^{n \times p} \\
\beta &\in \mathbb{R}^p \\
\epsilon &\in \mathbb{R}^n \\
\end{align*}

\end{frame}

%%%%%%%%%%%%%%%%%%%%%%%%%%%%%%%%%%%%%%%%%%%%%%%%%%%
\begin{frame}[fragile] \frametitle{}

{\bf Vector Norms}

When working with vectors and matricies, it will be helpful
to represent certain quantities by norms. The p-norm of a
vector is given by:
\begin{align*}
||x||_p^p = \sum_{i=1}^n | x_i |^p
\end{align*}
\pause In particular, the squared $2$-norm yields the sum
of squares of a vector.

\end{frame}

%%%%%%%%%%%%%%%%%%%%%%%%%%%%%%%%%%%%%%%%%%%%%%%%%%%
\begin{frame}[fragile] \frametitle{}

{\bf Vector Norm Properties}

The following properties are true of all vector norms,
for a scalar $\alpha$ and vectors $v_1$ and $v_2$.
\begin{align*}
|| \alpha v_1 || = |\alpha| \cdot || v_1 || \\
|| v_1 + v_2 || \leq || v_1 || + || v_2 || \\
\end{align*}

\end{frame}

%%%%%%%%%%%%%%%%%%%%%%%%%%%%%%%%%%%%%%%%%%%%%%%%%%%
\begin{frame}[fragile] \frametitle{}

{\bf p-Norm Properties}

$p$-norms have several additional properties that we will find
useful.

\pause Define $q$ such that:
\begin{align*}
\frac{1}{p} + \frac{1}{q} = 1
\end{align*}
The $q$-norm and $p$-norm are then said to be {\it dual} to one another.

\pause Notice that the $2$-norm is dual to itself.

\end{frame}

%%%%%%%%%%%%%%%%%%%%%%%%%%%%%%%%%%%%%%%%%%%%%%%%%%%
\begin{frame}[fragile] \frametitle{}

{\bf p-Norm Properties, cont.}

Hölder's inequality then yields
\begin{align*}
|v_1^t v_2 | \leq ||v_1||_p ||v_2||_q
\end{align*}
\pause As a special case, the Cauchy–Schwarz inequality
gives that:
\begin{align*}
|v_1^t v_2 |^2 \leq ||v_1||_2^2 ||v_2||_2^2
\end{align*}

\end{frame}


%%%%%%%%%%%%%%%%%%%%%%%%%%%%%%%%%%%%%%%%%%%%%%%%%%%
\begin{frame}[fragile] \frametitle{}

{\bf p-Norm Properties, cont.}

Finally, and of most importance for us today, note that
the squared $2$-norm is exactly equal to the self inner
product:
\begin{align*}
|| v_1 ||_2^2 = v_1^t v_1
\end{align*}

\end{frame}

%%%%%%%%%%%%%%%%%%%%%%%%%%%%%%%%%%%%%%%%%%%%%%%%%%%
\begin{frame}[fragile] \frametitle{}

{\bf Least squares (again)}

To estimate the least squares solution, which is again the
MLE for independent normal errors, we see that:
\begin{align*}
\widehat{\beta} \in \argmin_{b \in \mathbb{R}^p} \left\{ ||Y - X \beta ||_2^2 \right\}
\end{align*}
\pause Now using vector norms to denote the sum of
squares.

\end{frame}

%%%%%%%%%%%%%%%%%%%%%%%%%%%%%%%%%%%%%%%%%%%%%%%%%%%
\begin{frame}[fragile] \frametitle{}

It will be helpful to re-write the sum of squares as:
\begin{eqnarray*}
||Y - X \beta ||_2^2 &=& (Y - X\beta)^t (Y - X\beta) \\ \pause
&=& (Y^t - \beta^t X^t) (Y - X\beta) \\ \pause
&=& Y^tY - Y^t X\beta - \beta^t X^t Y + \beta^t X^t X\beta \\ \pause
&=& Y^tY - 2 Y^t X\beta + \beta^t X^t X\beta
\end{eqnarray*}

\end{frame}

%%%%%%%%%%%%%%%%%%%%%%%%%%%%%%%%%%%%%%%%%%%%%%%%%%%
\begin{frame}[fragile] \frametitle{}

{\bf Normal Equations}

In order to find the minimum of the sum of squares, we take the gradient
with respect to $\beta$ and set it equal to zero.

Recall that, for a vector $a$ and symmetric matrix $A$ :
\begin{align*}
\nabla_\beta a^t \beta = a \\
\nabla_\beta \beta^t A \beta = 2 A \beta
\end{align*}
\pause This gives the gradient of the sum of squares as:
\begin{align*}
\nabla_\beta ||Y - X \beta ||_2^2 &= \nabla_\beta \left(Y^tY - {\color{solarized@blue} 2 Y^t X\beta} + {\color{solarized@magenta} \beta^t X^t X\beta} \right)\\
&= {\color{solarized@magenta}2 X^t X \beta} - {\color{solarized@blue} 2 X^t y}
\end{align*}

\end{frame}

%%%%%%%%%%%%%%%%%%%%%%%%%%%%%%%%%%%%%%%%%%%%%%%%%%%
\begin{frame}[fragile] \frametitle{}

Setting this equal to zero gives a set of $p$ equations called
the normal equations:
\begin{align*}
X^t X \widehat{\beta} &= X^t y
\end{align*}

\end{frame}

%%%%%%%%%%%%%%%%%%%%%%%%%%%%%%%%%%%%%%%%%%%%%%%%%%%
\begin{frame}[fragile] \frametitle{}

{\bf Maximum or Minimum?}

To determine whether the normal equations give a local minimum, maximum, or
saddle point, we can calculate the Hessian matrix.
\pause This is a $p \times p$ matrix giving every combination of the
second partial derivatives:
\begin{align*}
H f(\beta) &=
  \left(\begin{array}{cccc}\frac{\partial^2f}{\partial \beta_1 \partial \beta_1}&\frac{\partial^2f}{\partial \beta_1 \partial \beta_2}&\cdots&\frac{\partial^2f}{\partial \beta_1 \partial \beta_p}\\
                           \frac{\partial^2f}{\partial \beta_2 \partial \beta_1}&\ddots&&\frac{\partial^2f}{\partial \beta_2 \partial \beta_p}\\
                           \vdots&&\ddots&\vdots\\
                           \frac{\partial^2f}{\partial \beta_n \partial \beta_1}&\frac{\partial^2f}{\partial \beta_n \partial \beta_2}&\cdots&\frac{\partial^2f}{\partial \beta_n \partial \beta_p}\\\end{array}\right)
\end{align*}
If the Hessian is positive definite ($x^t H x \geq 0$) at a critical point,
then the critical point is a local minimum.

\end{frame}

%%%%%%%%%%%%%%%%%%%%%%%%%%%%%%%%%%%%%%%%%%%%%%%%%%%
\begin{frame}[fragile] \frametitle{}

Looking at the gradiant of the sum of squares:
\begin{align*}
\nabla_\beta ||Y - X \beta ||_2^2 &= 2 X^t X \beta - 2 X^t y
\end{align*}
\pause We can see that the Hessian is simply:
\begin{align*}
H_\beta ||Y - X \beta ||_2^2 &= 2 X^t X
\end{align*}
\pause Why is this positive definite? \pause
\begin{align*}
v^t \left(2 X^tX \right) v &= 2 \left( v^t X^t X v\right) \\
&= 2 || X v ||_2^2 \\
&\geq 0
\end{align*}

\end{frame}

%%%%%%%%%%%%%%%%%%%%%%%%%%%%%%%%%%%%%%%%%%%%%%%%%%%
\begin{frame}[fragile] \frametitle{}

Back to the normal equations themselves, notice that if the
matrix $X^t X$ is invertable, we can `solve' the normal
equations as:
\begin{align*}
X^t X \widehat{\beta} &= X^t y \\
\widehat{\beta} &= (X^t X)^{-1} X^t y
\end{align*}
\pause This is not a good way to solve the normal equations
numerically, but for deriving theoretical results about the
least squares estimator this form will be very useful.

\end{frame}

%%%%%%%%%%%%%%%%%%%%%%%%%%%%%%%%%%%%%%%%%%%%%%%%%%%
\begin{frame}[fragile] \frametitle{}

\begin{flushright}
{\color{yaleblue}\sc\fontsize{1cm}{0cm}\selectfont Matricies and Model Frames in R}
\end{flushright}

\end{frame}


\end{document}











